\documentclass[12pt]{article}
\usepackage{latexblog}

% ============================================================
% Blog Metadata
% ============================================================
\blogtitle{在虚拟机中使用Docker}
\blogdate{2017-06-17}
\blogtags{Docker, 云计算, 容器化}
\bloglang{zh}
\blogsummary{}

\begin{document}
\makeblogtitle

几年前在公司复杂维护机房的时候，就开始关注Docker这种基于容器的虚拟化技术了，当时并没有选择Docker，因为几年前docker刚起步还不是很成熟，不敢采用这样的技术（当然关键是自己不了解也能力不够）。当时采取的是KVM和Virtual Box，问题也很明显，因为一台物理机（Dell T320 32GRAM)开个四五台Virtual Box虚拟机就有点吃不消了，想做到专机专用，也是很困难的事情。当时的主要目的是想把Oracle、WebSphere等吃内存的东西隔离出来。

现在有机会接触到Docker了，有必要认真的学习下了。貌似在Mac上Docker实际上是运行在虚拟的Linux中的，因此决定使用虚拟机来运行Docker，以下是我的配置：

\begin{itemize}
\tightlist
\item
  Mac Book Pro, OSX
\item
  Virtual Box, Ubuntu-server 16.04.2 X64 LTS, 4G Ram, 30G HDD，网卡桥接
\end{itemize}

好了，首先是安装Docker，参考\href{https://docs.docker.com/engine/installation/linux/ubuntu/\#install-docker}{官方文档}

\begin{Shaded}
\begin{Highlighting}[]
\FunctionTok{sudo}\NormalTok{ apt{-}get install }\DataTypeTok{\textbackslash{}}
\NormalTok{    apt{-}transport{-}https }\DataTypeTok{\textbackslash{}}
\NormalTok{    ca{-}certificates }\DataTypeTok{\textbackslash{}}
\NormalTok{    curl }\DataTypeTok{\textbackslash{}}
\NormalTok{    software{-}properties{-}common}
\ExtensionTok{curl} \AttributeTok{{-}fsSL}\NormalTok{ https://download.docker.com/linux/ubuntu/gpg }\KeywordTok{|} \FunctionTok{sudo}\NormalTok{ apt{-}key add }\AttributeTok{{-}}
\FunctionTok{sudo}\NormalTok{ add{-}apt{-}repository }\DataTypeTok{\textbackslash{}}
   \StringTok{"deb [arch=amd64] https://download.docker.com/linux/ubuntu }\DataTypeTok{\textbackslash{}}
\StringTok{   }\VariableTok{$(}\ExtensionTok{lsb\_release} \AttributeTok{{-}cs}\VariableTok{)}\StringTok{ }\DataTypeTok{\textbackslash{}}
\StringTok{   stable"}
\FunctionTok{sudo}\NormalTok{ apt{-}get update}
\FunctionTok{sudo}\NormalTok{ apt{-}get install docker{-}ce}
\end{Highlighting}
\end{Shaded}

好了，这样就安装成功了。来下载一些\href{https://hub.docker.com/}{docker镜像}吧，我们在后面可能慢慢会用到（先这样设想吧）：

\begin{itemize}
\tightlist
\item
  oracle/openjdk:8 选用Oracle的JDK8镜像来部署SpringBoot应用
\item
  gitlab/gitlab-ce 用来搭建私有的GIT仓库
\item
  sonatype/nexus3 用来搭建私有的Docker镜像库
\item
  gocd/gocd-server 用来搭建gocd作CICD
\item
  node 用来运行NodeJS的前端
\item
  percona 用来提供MySQL服务
\item
  mongo 用来提供MongoDB服务
\end{itemize}

\begin{Shaded}
\begin{Highlighting}[]
\ExtensionTok{docker}\NormalTok{ pull oracle/openjdk}
\ExtensionTok{docker}\NormalTok{ pull gitlab/gitlab{-}ce}
\ExtensionTok{docker}\NormalTok{ pull sonatype/nexus3}
\ExtensionTok{docker}\NormalTok{ pull gocd/gocd{-}server}
\ExtensionTok{docker}\NormalTok{ pull node}
\ExtensionTok{docker}\NormalTok{ pull percona}
\ExtensionTok{docker}\NormalTok{ pull mongo}
\end{Highlighting}
\end{Shaded}

你会发现太慢了，是不是?幸好可以有加速的方式，可以试用\href{https://www.daocloud.io/mirror\#accelerator-doc}{DaoCloud}的加速器。

\begin{Shaded}
\begin{Highlighting}[]
\FunctionTok{sudo}\NormalTok{ vim /etc/docker/daemon.json}
\end{Highlighting}
\end{Shaded}

输入以下内容:

\begin{Shaded}
\begin{Highlighting}[]
\FunctionTok{\{}
    \DataTypeTok{"registry{-}mirrors"}\FunctionTok{:} \OtherTok{[}
        \StringTok{"http://1729****.m.daocloud.io"}
    \OtherTok{]}\FunctionTok{,}
    \DataTypeTok{"insecure{-}registries"}\FunctionTok{:} \OtherTok{[]}
\FunctionTok{\}}
\end{Highlighting}
\end{Shaded}

重启Docker：

\begin{Shaded}
\begin{Highlighting}[]
\FunctionTok{sudo}\NormalTok{ /etc/init.d/docker restart}
\end{Highlighting}
\end{Shaded}


\end{document}
